
% Dateiname: MoE_Glossar_Apeiron_261029-1.tex

\documentclass[11pt,a4paper]{article}

\usepackage[T1]{fontenc}
\usepackage[utf8]{inputenc}
\usepackage[ngerman]{babel}
\usepackage[a4paper,margin=25mm]{geometry}
\usepackage{amsmath,amssymb}
\usepackage{helvet}
\usepackage{sfmath}
\renewcommand{\familydefault}{\sfdefault}
\usepackage{microtype}
\usepackage[hidelinks]{hyperref}
\usepackage{lastpage}
\usepackage{fancyhdr}
\usepackage{parskip}

\pagestyle{fancy}
\fancyhf{}
\fancyfoot[C]{\thepage\ (\pageref{LastPage})}
\renewcommand{\headrulewidth}{0pt}
\renewcommand{\footrulewidth}{0pt}

\hypersetup{
  pdftitle={MoE-Glossar: Apeiron und Apeiron-Expansion},
  pdfauthor={}
}

\begin{document}

\vspace*{-1.8cm}

\noindent
\begin{minipage}[t]{0.66\textwidth}
{\LARGE\bfseries Apeiron und
Apeiron-Expansion}
\end{minipage}
\hfill
\begin{minipage}[t]{0.30\textwidth}
\raggedleft
{\small MoE-Glossareintrag}\\
{\small Version 261029-1}
\end{minipage}

\vspace{1.4em}

\section*{Apeiron}

Das Apeiron ist die im Model of
Everything (MoE) zugrunde gelegte
kontinuierliche Ausgangsgröße.

Seine räumliche Verteilung wird
durch die Apeiron-Dichte
$\rho_A$ beschrieben.

Die in einem betrachteten Volumen
$V_A$ enthaltene Apeiron-Menge
wird mit $A$ bezeichnet.

\section*{Apeiron-Expansion}

Die Apeiron-Expansion ist der
grundlegende Prozess der
räumlichen Volumenzunahme.

Mit der Expansion verändert
sich die Apeiron-Dichte.

Für eine erhaltene Apeiron-Menge
gilt:

\[
A=\rho_A V_A=\mathrm{konstant}.
\]

Daraus folgt:

\[
\rho_A=\frac{A}{V_A}.
\]

Die Änderung der Dichte mit
dem Volumen lautet:

\[
\frac{\mathrm{d}\rho_A}
{\mathrm{d}V_A}
=
-\frac{A}{V_A^2}.
\]

Die Expansion bewirkt somit
eine Verdünnung des Apeirons.

\section*{Normierung}

Für die normierte Apeiron-Menge
$A=1$ gilt:

\[
\rho_A V_A=1.
\]

Der Normpunkt ist:

\[
(V_A,\rho_A)=(1,1).
\]

\section*{Prozessparameter}

Der Verlauf der Expansion
wird durch den Prozessparameter
$\tau$ beschrieben.

Volumen und Dichte sind
Funktionen dieses Parameters:

\[
V_A=V_A(\tau),
\qquad
\rho_A=\rho_A(\tau).
\]

Die Erhaltungsbedingung lautet:

\[
\rho_A(\tau)V_A(\tau)=A.
\]

Durch Differentiation nach
$\tau$ folgt:

\[
\frac{\mathrm{d}\rho_A}
{\mathrm{d}\tau}V_A
+
\rho_A
\frac{\mathrm{d}V_A}
{\mathrm{d}\tau}
=0.
\]

Damit ergibt sich:

\[
\frac{\mathrm{d}\rho_A}
{\mathrm{d}\tau}
=
-\frac{\rho_A}{V_A}
\frac{\mathrm{d}V_A}
{\mathrm{d}\tau}.
\]

Der Prozessparameter $\tau$
ordnet die aufeinanderfolgenden
Zustände der Expansion.

\section*{Zeit}

Zeit bezeichnet im MoE den
Fortgang des Expansionsprozesses.

Die zeitliche Ordnung physikalischer
Vorgänge wird durch die Abfolge
der Expansionszustände bestimmt.

Der Prozessparameter $\tau$
dient ihrer mathematischen
Beschreibung.

\section*{Lokale Expansion}

Die räumliche Verteilung
des Apeirons bestimmt die
lokalen Expansionsbedingungen.

Die lokale Expansionsgeschwindigkeit
wird im MoE mit $c$ bezeichnet.

Dichteunterschiede und
Expansionsunterschiede stehen
in einem kausalen Zusammenhang.

\section*{Quantelung und Raumbildung}

Innerhalb des expandierenden
Apeirons entstehen lokale
Dichte- und
Geschwindigkeitsunterschiede.

Dynamische Scherung führt
zur Ausbildung von Rotation.

Durch Rotationsschließung
entstehen QUBLES.

Ihre räumlichen Kontaktordnungen
führen zur Ausbildung der
Superdichten QUBLE-Packung (SDQP)
und der Dichten QUBLE-Packung (DQP).

\section*{Emergenzfolge}

\[
\begin{gathered}
\text{Apeiron}
\rightarrow
\text{Expansion}
\rightarrow
\text{Quantelung}
\\[0.4em]
\rightarrow
\text{QUBLE}
\rightarrow
\text{SDQP}
\rightarrow
\text{DQP}
\rightarrow
\text{Pfad}
\end{gathered}
\]

\vspace{0.9em}
\hrule
\vspace{0.45em}

{\footnotesize
MoE-Glossar:
\emph{QUBLE / Quantelung},
\emph{SDQP},
\emph{DQP},
\emph{Interstitium},
\emph{Pfad}.
}

\end{document}
